% Options for packages loaded elsewhere
\PassOptionsToPackage{unicode}{hyperref}
\PassOptionsToPackage{hyphens}{url}
\documentclass[
]{book}
\usepackage{xcolor}
\usepackage{amsmath,amssymb}
\setcounter{secnumdepth}{5}
\usepackage{iftex}
\ifPDFTeX
  \usepackage[T1]{fontenc}
  \usepackage[utf8]{inputenc}
  \usepackage{textcomp} % provide euro and other symbols
\else % if luatex or xetex
  \usepackage{unicode-math} % this also loads fontspec
  \defaultfontfeatures{Scale=MatchLowercase}
  \defaultfontfeatures[\rmfamily]{Ligatures=TeX,Scale=1}
\fi
\usepackage{lmodern}
\ifPDFTeX\else
  % xetex/luatex font selection
\fi
% Use upquote if available, for straight quotes in verbatim environments
\IfFileExists{upquote.sty}{\usepackage{upquote}}{}
\IfFileExists{microtype.sty}{% use microtype if available
  \usepackage[]{microtype}
  \UseMicrotypeSet[protrusion]{basicmath} % disable protrusion for tt fonts
}{}
\makeatletter
\@ifundefined{KOMAClassName}{% if non-KOMA class
  \IfFileExists{parskip.sty}{%
    \usepackage{parskip}
  }{% else
    \setlength{\parindent}{0pt}
    \setlength{\parskip}{6pt plus 2pt minus 1pt}}
}{% if KOMA class
  \KOMAoptions{parskip=half}}
\makeatother
\usepackage{color}
\usepackage{fancyvrb}
\newcommand{\VerbBar}{|}
\newcommand{\VERB}{\Verb[commandchars=\\\{\}]}
\DefineVerbatimEnvironment{Highlighting}{Verbatim}{commandchars=\\\{\}}
% Add ',fontsize=\small' for more characters per line
\usepackage{framed}
\definecolor{shadecolor}{RGB}{248,248,248}
\newenvironment{Shaded}{\begin{snugshade}}{\end{snugshade}}
\newcommand{\AlertTok}[1]{\textcolor[rgb]{0.94,0.16,0.16}{#1}}
\newcommand{\AnnotationTok}[1]{\textcolor[rgb]{0.56,0.35,0.01}{\textbf{\textit{#1}}}}
\newcommand{\AttributeTok}[1]{\textcolor[rgb]{0.13,0.29,0.53}{#1}}
\newcommand{\BaseNTok}[1]{\textcolor[rgb]{0.00,0.00,0.81}{#1}}
\newcommand{\BuiltInTok}[1]{#1}
\newcommand{\CharTok}[1]{\textcolor[rgb]{0.31,0.60,0.02}{#1}}
\newcommand{\CommentTok}[1]{\textcolor[rgb]{0.56,0.35,0.01}{\textit{#1}}}
\newcommand{\CommentVarTok}[1]{\textcolor[rgb]{0.56,0.35,0.01}{\textbf{\textit{#1}}}}
\newcommand{\ConstantTok}[1]{\textcolor[rgb]{0.56,0.35,0.01}{#1}}
\newcommand{\ControlFlowTok}[1]{\textcolor[rgb]{0.13,0.29,0.53}{\textbf{#1}}}
\newcommand{\DataTypeTok}[1]{\textcolor[rgb]{0.13,0.29,0.53}{#1}}
\newcommand{\DecValTok}[1]{\textcolor[rgb]{0.00,0.00,0.81}{#1}}
\newcommand{\DocumentationTok}[1]{\textcolor[rgb]{0.56,0.35,0.01}{\textbf{\textit{#1}}}}
\newcommand{\ErrorTok}[1]{\textcolor[rgb]{0.64,0.00,0.00}{\textbf{#1}}}
\newcommand{\ExtensionTok}[1]{#1}
\newcommand{\FloatTok}[1]{\textcolor[rgb]{0.00,0.00,0.81}{#1}}
\newcommand{\FunctionTok}[1]{\textcolor[rgb]{0.13,0.29,0.53}{\textbf{#1}}}
\newcommand{\ImportTok}[1]{#1}
\newcommand{\InformationTok}[1]{\textcolor[rgb]{0.56,0.35,0.01}{\textbf{\textit{#1}}}}
\newcommand{\KeywordTok}[1]{\textcolor[rgb]{0.13,0.29,0.53}{\textbf{#1}}}
\newcommand{\NormalTok}[1]{#1}
\newcommand{\OperatorTok}[1]{\textcolor[rgb]{0.81,0.36,0.00}{\textbf{#1}}}
\newcommand{\OtherTok}[1]{\textcolor[rgb]{0.56,0.35,0.01}{#1}}
\newcommand{\PreprocessorTok}[1]{\textcolor[rgb]{0.56,0.35,0.01}{\textit{#1}}}
\newcommand{\RegionMarkerTok}[1]{#1}
\newcommand{\SpecialCharTok}[1]{\textcolor[rgb]{0.81,0.36,0.00}{\textbf{#1}}}
\newcommand{\SpecialStringTok}[1]{\textcolor[rgb]{0.31,0.60,0.02}{#1}}
\newcommand{\StringTok}[1]{\textcolor[rgb]{0.31,0.60,0.02}{#1}}
\newcommand{\VariableTok}[1]{\textcolor[rgb]{0.00,0.00,0.00}{#1}}
\newcommand{\VerbatimStringTok}[1]{\textcolor[rgb]{0.31,0.60,0.02}{#1}}
\newcommand{\WarningTok}[1]{\textcolor[rgb]{0.56,0.35,0.01}{\textbf{\textit{#1}}}}
\usepackage{longtable,booktabs,array}
\usepackage{calc} % for calculating minipage widths
% Correct order of tables after \paragraph or \subparagraph
\usepackage{etoolbox}
\makeatletter
\patchcmd\longtable{\par}{\if@noskipsec\mbox{}\fi\par}{}{}
\makeatother
% Allow footnotes in longtable head/foot
\IfFileExists{footnotehyper.sty}{\usepackage{footnotehyper}}{\usepackage{footnote}}
\makesavenoteenv{longtable}
\usepackage{graphicx}
\makeatletter
\newsavebox\pandoc@box
\newcommand*\pandocbounded[1]{% scales image to fit in text height/width
  \sbox\pandoc@box{#1}%
  \Gscale@div\@tempa{\textheight}{\dimexpr\ht\pandoc@box+\dp\pandoc@box\relax}%
  \Gscale@div\@tempb{\linewidth}{\wd\pandoc@box}%
  \ifdim\@tempb\p@<\@tempa\p@\let\@tempa\@tempb\fi% select the smaller of both
  \ifdim\@tempa\p@<\p@\scalebox{\@tempa}{\usebox\pandoc@box}%
  \else\usebox{\pandoc@box}%
  \fi%
}
% Set default figure placement to htbp
\def\fps@figure{htbp}
\makeatother
\setlength{\emergencystretch}{3em} % prevent overfull lines
\providecommand{\tightlist}{%
  \setlength{\itemsep}{0pt}\setlength{\parskip}{0pt}}
\usepackage[]{natbib}
\bibliographystyle{apalike}
\usepackage{bookmark}
\IfFileExists{xurl.sty}{\usepackage{xurl}}{} % add URL line breaks if available
\urlstyle{same}
\hypersetup{
  pdftitle={Endocrine Networks Simulations},
  pdfauthor={Dr Massoud Boroujerdi},
  hidelinks,
  pdfcreator={LaTeX via pandoc}}

\title{Endocrine Networks Simulations}
\author{Dr Massoud Boroujerdi}
\date{2026-02-24}

\begin{document}
\maketitle

\chapter{Methods}\label{Methods}

\section{Mathematical Framework}\label{mathematical-framework}

\subsection{Second-Order Negative Feedback Model}\label{second-order-negative-feedback-model}

The fundamental building block of our network model is the second-order negative feedback compartmental system \citep{Boroujerdi2013, Boroujerdi2024}. Each element (substrate or hormone) is represented by two compartments: S (substrate/hormone) and C (controller), governed by the differential equations:

\[\frac{dS}{dt} = R_a - k_3C\]

\[\frac{dC}{dt} = k_4S - k_2C\]

At steady state, the production rate relates to concentration via the metabolic clearance rate (MCR):

\[R_a = \frac{k_3k_4}{k_2}S = \text{MCR} \cdot S\]

where \(\text{MCR} = \frac{k_3k_4}{k_2}\) with dimensions of time\(^{-1}\).

\subsection{Transfer Function Formulation}\label{transfer-function-formulation}

The system parameters relate to the standard second-order differential equation parameters:
- Natural frequency: \(\omega_n = \sqrt{k_3k_4}\)
- Damping ratio: \(\zeta = \frac{k_2}{2\omega_n}\)

Equivalently: \(k_2 = 2\zeta\omega_n\) and \(k_3 = k_4 = \omega_n\).

The state-space representation is:

\[\dot{\mathbf{x}} = A\mathbf{x} + Bu\]

\[y = C\mathbf{x}\]

where:

\[A = \begin{bmatrix} 0 & -k_3 \\ k_4 & -k_2 \end{bmatrix}, \quad B = \begin{bmatrix} 1 \\ 0 \end{bmatrix}\]

\subsection{G-Type and H-Type Transfer Functions}\label{g-type-and-h-type-transfer-functions}

Two transfer function architectures are considered, corresponding to different controller flux connections:

\textbf{G-type (feed-forward):} Uses the \(k_3C\) flux, yielding:

\[G_i(s) = \frac{k_{3i}(s + k_{2i})}{s(s + k_{2i}) + k_{3i}k_{4i}}\]

This transfer function includes a zero at \(s = -k_{2i}\), providing phase lead compensation and faster response.

\textbf{H-type (feedback):} Uses the \(k_2C\) flux, yielding:

\[H_i(s) = \frac{k_{2i}k_{4i}}{s(s + k_{2i}) + k_{3i}k_{4i}}\]

This transfer function has no zeros, resulting in pure lag behavior with enhanced damping and stronger noise rejection.

\subsection{Network Composition}\label{network-composition}

For a network of \(N\) elements, the overall transfer function is obtained by cascade multiplication:

\[G_{\text{network}}(s) = \prod_{i=1}^{N} T_i(s)\]

where \(T_i(s)\) is either \(G_i(s)\) or \(H_i(s)\) depending on the architecture assignment.

For networks with feedback loops, the closed-loop transfer function is:

\[G_{\text{closed}}(s) = \frac{G_{\text{forward}}(s)}{1 + G_{\text{forward}}(s)G_{\text{feedback}}(s)}\]

\subsection{Parameter Assignment from MCR}\label{parameter-assignment-from-mcr}

For each element \(i\) representing a substrate or hormone:

\begin{enumerate}
\def\labelenumi{\arabic{enumi}.}
\item
  \textbf{Obtain MCR:} From literature or experimental measurements (units: time\(^{-1}\))
\item
  \textbf{Compute natural frequency:}
  \[\omega_{n,i} = \text{MCR}_i \cdot \text{scaling factor}\]

  The scaling factor relates MCR to the characteristic frequency of the system. For simplicity, we use \(\omega_{n,i} = \text{MCR}_i\) as a first approximation.
\item
  \textbf{Assign damping ratio:}

  \begin{itemize}
  \tightlist
  \item
    For critically damped systems: \(\zeta_i = 1.0\)
  \item
    For under damped systems (potential oscillations): \(\zeta_i = 0.5\) to \(0.7\)
  \item
    From experimental data when available
  \end{itemize}
\item
  \textbf{Calculate rate constants:}
  \[k_{2i} = 2\zeta_i\omega_{n,i}\]
  \[k_{3i} = k_{4i} = \omega_{n,i}\]
\end{enumerate}

\subsection{Time-Domain Solutions via Inverse Laplace Transform}\label{time-domain-solutions-via-inverse-laplace-transform}

For step input \(u(t) = u_0 \cdot \mathbb{1}(t)\) (unit step function), the output in the time domain is obtained via inverse Laplace transform:

\[y(t) = \mathcal{L}^{-1}\left\{G(s) \cdot \frac{u_0}{s}\right\}\]

For a second-order system with transfer function:

\[G(s) = \frac{\omega_n^2}{s^2 + 2\zeta\omega_n s + \omega_n^2}\]

The analytical solution for a unit step input depends on damping:

\textbf{Underdamped case} (\(\zeta < 1\)):

\[y(t) = 1 - \frac{e^{-\zeta\omega_n t}}{\sqrt{1-\zeta^2}} \sin\left(\omega_d t + \phi\right)\]

where \(\omega_d = \omega_n\sqrt{1-\zeta^2}\) (damped natural frequency) and \(\phi = \arctan\left(\frac{\sqrt{1-\zeta^2}}{\zeta}\right)\).

\textbf{Critically damped case} (\(\zeta = 1\)):

\[y(t) = 1 - e^{-\omega_n t}(1 + \omega_n t)\]

\textbf{Overdamped case} (\(\zeta > 1\)):

\[y(t) = 1 - \frac{1}{2\omega_n\sqrt{\zeta^2-1}}\left[(\zeta + \sqrt{\zeta^2-1})e^{s_1 t} - (\zeta - \sqrt{\zeta^2-1})e^{s_2 t}\right]\]

where \(s_1 = -\omega_n(\zeta + \sqrt{\zeta^2-1})\) and \(s_2 = -\omega_n(\zeta - \sqrt{\zeta^2-1})\).

\section{Frequency-Domain Analysis}\label{frequency-domain-analysis}

\subsection{Bode Plot Construction}\label{bode-plot-construction}

Frequency response is evaluated by substituting \(s = j\omega\) into the transfer function:

\[H(j\omega) = G(j\omega)\]

The magnitude (in dB) and phase are:

\[|H(j\omega)|_{\text{dB}} = 20\log_{10}|H(j\omega)|\]

\[\angle H(j\omega) = \arctan\left(\frac{\text{Im}[H(j\omega)]}{\text{Re}[H(j\omega)]}\right)\]

Bode plots are generated by evaluating \(H(j\omega)\) across a logarithmic frequency range (typically \(10^{-3}\) to \(10^{3}\) rad/time-unit).

\subsection{Performance Metrics}\label{performance-metrics}

From the Bode plot, we extract:

\begin{enumerate}
\def\labelenumi{\arabic{enumi}.}
\tightlist
\item
  \textbf{DC gain:} \(|H(j\omega)|_{\text{dB}}\) at \(\omega \to 0\)
\item
  \textbf{Cutoff frequency (\(\omega_c\)):} Frequency where \(|H(j\omega)|_{\text{dB}} = |H(0)|_{\text{dB}} - 3\) dB
\item
  \textbf{Bandwidth:} Range from DC to \(\omega_c\)
\item
  \textbf{High-frequency attenuation:} \(|H(j\omega)|_{\text{dB}}\) at \(\omega \to \infty\)
\item
  \textbf{Phase margin:} Phase at gain crossover frequency
\item
  \textbf{Gain margin:} Gain at phase crossover frequency (\(-180°\))
\end{enumerate}

These metrics directly translate to biological properties:
- \textbf{Cutoff frequency} → characteristic response time (\(\tau \approx 2\pi/\omega_c\))
- \textbf{Bandwidth} → ability to track time-varying signals
- \textbf{High-frequency attenuation} → noise rejection capability
- \textbf{Phase margin} → stability reserve

\subsection{Stability Analysis}\label{stability-analysis}

Network stability is determined by the locations of the poles of \(G_{\text{network}}(s)\):

\textbf{Routh-Hurwitz criterion:} System is stable if all coefficients of the characteristic polynomial are positive and all elements in the first column of the Routh array are positive.

\textbf{Eigenvalue criterion:} For state-space representation, compute eigenvalues \(\lambda\) of matrix \(A\). System is stable if \(\text{Re}(\lambda_i) < 0\) for all \(i\).

\textbf{Nyquist criterion:} System is stable if the Nyquist plot of the open-loop transfer function \(L(j\omega) = G_{\text{forward}}(j\omega)G_{\text{feedback}}(j\omega)\) does not encircle the point \((-1, 0)\).

\section{Computational Implementation}\label{computational-implementation}

\subsection{Software and Packages}\label{software-and-packages}

All simulations were performed in R (version 4.3.0 or higher). Key packages:

\begin{itemize}
\tightlist
\item
  \texttt{deSolve}: Numerical solution of differential equations \citep{Soetaert2010}
\item
  \texttt{control}: Control systems analysis (if available; alternatively custom implementation)
\item
  \texttt{signal}: Signal processing and frequency analysis \citep{Signal}
\item
  \texttt{ggplot2}: Data visualization \citep{Wickham2016}
\item
  \texttt{bookdown}: Document preparation \citep{Xie2016}
\end{itemize}

The algorithms implementing the network transfer function, frequency response, and time-domain simulation are provided in the \hyperref[Appendix]{Appendix}.

\section{Validation Approach}\label{validation-approach}

\subsection{Comparison with Known Systems}\label{comparison-with-known-systems}

Model predictions are validated against experimental data from well-characterized systems:

\begin{enumerate}
\def\labelenumi{\arabic{enumi}.}
\tightlist
\item
  \textbf{Glucose-insulin system:}

  \begin{itemize}
  \tightlist
  \item
    Oral glucose tolerance test (OGTT) data
  \item
    Ultradian oscillation period (\textasciitilde90 min)
  \item
    Insulin clearance rate (\textasciitilde5 min half-life)
  \end{itemize}
\item
  \textbf{HPA axis (cortisol):}

  \begin{itemize}
  \tightlist
  \item
    Cortisol circadian rhythm
  \item
    Response to ACTH stimulation
  \item
    Dexamethasone suppression dynamics
  \end{itemize}
\item
  \textbf{Thyroid axis:}

  \begin{itemize}
  \tightlist
  \item
    TSH response to TRH stimulation
  \item
    T4 and T3 half-lives (7 days and 1 day, respectively)
  \item
    Levothyroxine dose-response curves
  \end{itemize}
\end{enumerate}

\subsection{Sensitivity Analysis}\label{sensitivity-analysis}

Parameter sensitivity is assessed by:

\begin{enumerate}
\def\labelenumi{\arabic{enumi}.}
\tightlist
\item
  \textbf{One-at-a-time (OAT):} Vary each parameter individually by ±20\% while holding others constant
\item
  \textbf{Monte Carlo:} Random sampling from parameter distributions, compute output statistics
\item
  \textbf{Sobol indices:} Variance-based sensitivity analysis to quantify parameter importance
\end{enumerate}

\section{Data Sources}\label{data-sources}

\subsection{Metabolic Clearance Rates}\label{metabolic-clearance-rates}

MCR values are obtained from published literature and pharmacokinetic databases:

\begin{itemize}
\tightlist
\item
  \textbf{Glucose:} Measured via glucose clamp studies \citep{DeFronzo1979}
\item
  \textbf{Insulin:} From C-peptide kinetics or insulin infusion studies \citep{Polonsky1986}
\item
  \textbf{Cortisol:} From tracer dilution studies \citep{Rosenfield1971}
\item
  \textbf{Thyroid hormones:} From radiotracer kinetics \citep{Cavalieri1984}
\item
  \textbf{Growth hormone:} From pulsatile secretion studies \citep{Veldhuis1996}
\end{itemize}

\subsection{Network Topologies}\label{network-topologies}

Feedback connections are determined from established physiological knowledge:

\begin{itemize}
\tightlist
\item
  Endocrine axes: Hypothalamic-pituitary-peripheral organ structure
\item
  Glucose-insulin: Glucose stimulates insulin, insulin suppresses hepatic glucose output
\item
  Thyroid: Negative feedback of T3 on TRH and TSH
\end{itemize}

\section{Reporting Standards}\label{reporting-standards}

Results are reported following control systems engineering conventions:

\begin{itemize}
\tightlist
\item
  \textbf{Time constants:} Reported as \(\tau = 1/\omega_n\) or response time to 63\% of final value
\item
  \textbf{Settling time:} Time to reach and stay within 2\% of final value
\item
  \textbf{Rise time:} Time to go from 10\% to 90\% of final value
\item
  \textbf{Overshoot:} \((y_{\max} - y_{\infty}) / y_{\infty} \times 100\%\)
\end{itemize}

\chapter*{Example of values}\label{example-of-values}
\addcontentsline{toc}{chapter}{Example of values}

\begin{longtable}[t]{llrrr}
\caption{\label{tab:table-1}Example of values (Part 1)}\\
\toprule
name & category & MCR\_min & MCR\_h & omega\_n\_rad\_per\_min\\
\midrule
Glucose & substrate & 0.06670 & 4.0000 & 0.06670\\
Free\_Fatty\_Acids & substrate & 0.04620 & 2.7700 & 0.04620\\
Amino\_Acids\_Mixed & substrate & 0.02310 & 1.3900 & 0.02310\\
Lactate & substrate & 0.11550 & 6.9300 & 0.11550\\
Ketone\_Bodies & substrate & 0.03470 & 2.0800 & 0.03470\\
\addlinespace
Insulin & hormone & 0.13860 & 8.3200 & 0.13860\\
Glucagon & hormone & 0.08660 & 5.2000 & 0.08660\\
Growth\_Hormone & hormone & 0.03470 & 2.0800 & 0.03470\\
Cortisol & hormone & 0.00770 & 0.4620 & 0.00770\\
ACTH & hormone & 0.06930 & 4.1600 & 0.06930\\
\addlinespace
Epinephrine & hormone & 0.34650 & 20.8000 & 0.34650\\
Norepinephrine & hormone & 0.23100 & 13.9000 & 0.23100\\
Leptin & hormone & 0.00140 & 0.0830 & 0.00140\\
Ghrelin & hormone & 0.02310 & 1.3900 & 0.02310\\
IGF1 & hormone & 0.00030 & 0.0170 & 0.00030\\
\addlinespace
TSH & hormone & 0.01160 & 0.6940 & 0.01160\\
T4\_Thyroxine & hormone & 0.00007 & 0.0042 & 0.00007\\
T3\_Triiodothyronine & hormone & 0.00048 & 0.0290 & 0.00048\\
Testosterone & hormone & 0.00100 & 0.0600 & 0.00100\\
Estradiol & hormone & 0.03470 & 2.0800 & 0.03470\\
\addlinespace
Progesterone & hormone & 0.00200 & 0.1200 & 0.00200\\
LH & hormone & 0.02310 & 1.3900 & 0.02310\\
FSH & hormone & 0.00480 & 0.2880 & 0.00480\\
Prolactin & hormone & 0.03470 & 2.0800 & 0.03470\\
Oxytocin & hormone & 0.34650 & 20.8000 & 0.34650\\
\addlinespace
Vasopressin\_ADH & hormone & 0.06930 & 4.1600 & 0.06930\\
Parathyroid\_Hormone & hormone & 0.13860 & 8.3200 & 0.13860\\
Calcitonin & hormone & 0.06930 & 4.1600 & 0.06930\\
Vitamin\_D3\_Active & hormone & 0.00030 & 0.0170 & 0.00030\\
Renin & hormone & 0.03470 & 2.0800 & 0.03470\\
\addlinespace
Angiotensin\_II & hormone & 0.69300 & 41.6000 & 0.69300\\
Aldosterone & hormone & 0.03470 & 2.0800 & 0.03470\\
Melatonin & hormone & 0.03470 & 2.0800 & 0.03470\\
GLP1 & hormone & 0.34650 & 20.8000 & 0.34650\\
GIP & hormone & 0.11550 & 6.9300 & 0.11550\\
\addlinespace
Amylin & hormone & 0.13860 & 8.3200 & 0.13860\\
\bottomrule
\end{longtable}

\begin{longtable}[t]{rlrll}
\caption{\label{tab:table-2}Example of values continuation (Part 2):}\\
\toprule
zeta & suggested\_type & halflife\_min & reference\_key & notes\\
\midrule
0.7 & G & 10.4 & DeFronzo1979 & Glucose clamp studies; MCR \textasciitilde{}4.0 mL/kg/min\\
0.7 & G & 15.0 & Klein1986 & Tracer dilution; context-dependent\\
0.7 & G & 30.0 & Wolfe2006 & Varies by amino acid; leucine reference\\
0.7 & G & 6.0 & Stanley1985 & Rapid turnover\\
0.7 & G & 20.0 & Balasse1989 & Beta-hydroxybutyrate reference\\
\addlinespace
0.5 & G & 5.0 & Polonsky1986 & C-peptide kinetics; rapid clearance\\
0.5 & G & 8.0 & Lefebvre1975 & Pulsatile secretion\\
0.5 & H & 20.0 & Veldhuis1996 & Pulsatile secretion; GH pulses\\
0.7 & H & 90.0 & Rosenfield1971 & HPA axis; ultradian rhythms\\
0.5 & G & 10.0 & Orth1992 & Precursor to cortisol\\
\addlinespace
0.5 & G & 2.0 & Cryer1980 & Rapid clearance; stress response\\
0.5 & G & 3.0 & Cryer1980 & Rapid clearance; sympathetic\\
0.7 & H & 500.0 & Klein2004 & Slow clearance; adipose signal\\
0.5 & G & 30.0 & Akamizu2006 & Appetite regulation\\
0.7 & H & 2400.0 & Juul1994 & Very slow clearance; bound to BPs\\
\addlinespace
0.7 & H & 60.0 & Cavalieri1984 & Thyroid axis; pituitary\\
0.9 & H & 10080.0 & Cavalieri1984 & 7-day half-life; highly stable\\
0.8 & H & 1440.0 & Cavalieri1984 & 1-day half-life; active form\\
0.7 & H & 693.0 & Winters1969 & HPG axis; male\\
0.7 & H & 20.0 & Longcope1986 & HPG axis; female; cyclic\\
\addlinespace
0.7 & H & 347.0 & Little1972 & HPG axis; luteal phase\\
0.5 & G & 30.0 & Veldhuis1988 & Pulsatile; gonadotropin\\
0.7 & H & 144.0 & Veldhuis1988 & Gonadotropin; slower than LH\\
0.7 & H & 20.0 & Sassin1972 & Lactation; circadian rhythm\\
0.5 & G & 2.0 & Leake1981 & Pulsatile; rapid clearance\\
\addlinespace
0.5 & G & 10.0 & Robertson1976 & Osmotic regulation\\
0.5 & G & 5.0 & Zaheer2018 & Calcium homeostasis\\
0.5 & G & 10.0 & Hirsch1978 & Opposes PTH\\
0.9 & H & 2400.0 & Jones2014 & Calcidiol; very slow\\
0.5 & G & 20.0 & Nussberger1987 & RAAS; enzymatic\\
\addlinespace
0.5 & G & 1.0 & Campbell1990 & RAAS; vasoconstrictor\\
0.7 & H & 20.0 & Vecsei1979 & RAAS; mineralocorticoid\\
0.5 & G & 20.0 & Waldhauser1984 & Circadian; pineal\\
0.5 & G & 2.0 & Deacon2004 & Incretin; rapid DPP4 degradation\\
0.5 & G & 6.0 & Deacon2004 & Incretin; GI peptide\\
\addlinespace
0.5 & G & 5.0 & Young1993 & Co-secreted with insulin\\
\bottomrule
\end{longtable}

All simulations and analyses are reproducible via documented R scripts provided in supplementary materials.

\section{Appendix}\label{Appendix}

\textbf{R programs structure and code snippets}

\subsection{Algorithm: Network Transfer Function}\label{algorithm-network-transfer-function}

\begin{Shaded}
\begin{Highlighting}[]
\CommentTok{\# Pseudocode for network transfer function construction}

\NormalTok{network\_tf }\OtherTok{\textless{}{-}} \ControlFlowTok{function}\NormalTok{(elements, architecture) \{}
  \CommentTok{\# elements: list of (omega\_n, zeta, MCR) for each node}
  \CommentTok{\# architecture: vector of "G" or "H" for each node}
  
  \CommentTok{\# Initialize numerator and denominator}
\NormalTok{  num }\OtherTok{\textless{}{-}} \DecValTok{1}
\NormalTok{  den }\OtherTok{\textless{}{-}} \DecValTok{1}
  
  \ControlFlowTok{for}\NormalTok{ (i }\ControlFlowTok{in} \DecValTok{1}\SpecialCharTok{:}\FunctionTok{length}\NormalTok{(elements)) \{}
\NormalTok{    omega\_n }\OtherTok{\textless{}{-}}\NormalTok{ elements[[i]]}\SpecialCharTok{$}\NormalTok{omega\_n}
\NormalTok{    zeta }\OtherTok{\textless{}{-}}\NormalTok{ elements[[i]]}\SpecialCharTok{$}\NormalTok{zeta}
    
\NormalTok{    k2 }\OtherTok{\textless{}{-}} \DecValTok{2} \SpecialCharTok{*}\NormalTok{ zeta }\SpecialCharTok{*}\NormalTok{ omega\_n}
\NormalTok{    k3 }\OtherTok{\textless{}{-}}\NormalTok{ omega\_n}
\NormalTok{    k4 }\OtherTok{\textless{}{-}}\NormalTok{ omega\_n}
    
    \ControlFlowTok{if}\NormalTok{ (architecture[i] }\SpecialCharTok{==} \StringTok{"G"}\NormalTok{) \{}
      \CommentTok{\# G{-}type: numerator has (s + k2)}
\NormalTok{      num\_i }\OtherTok{\textless{}{-}} \FunctionTok{c}\NormalTok{(k3, k3}\SpecialCharTok{*}\NormalTok{k2)}
\NormalTok{      den\_i }\OtherTok{\textless{}{-}} \FunctionTok{c}\NormalTok{(}\DecValTok{1}\NormalTok{, k2, k3}\SpecialCharTok{*}\NormalTok{k4)}
\NormalTok{    \} }\ControlFlowTok{else}\NormalTok{ \{}
      \CommentTok{\# H{-}type: numerator is constant}
\NormalTok{      num\_i }\OtherTok{\textless{}{-}} \FunctionTok{c}\NormalTok{(k2}\SpecialCharTok{*}\NormalTok{k4)}
\NormalTok{      den\_i }\OtherTok{\textless{}{-}} \FunctionTok{c}\NormalTok{(}\DecValTok{1}\NormalTok{, k2, k3}\SpecialCharTok{*}\NormalTok{k4)}
\NormalTok{    \}}
    
    \CommentTok{\# Convolve (multiply polynomials)}
\NormalTok{    num }\OtherTok{\textless{}{-}} \FunctionTok{convolve}\NormalTok{(num, }\FunctionTok{rev}\NormalTok{(num\_i))}
\NormalTok{    den }\OtherTok{\textless{}{-}} \FunctionTok{convolve}\NormalTok{(den, }\FunctionTok{rev}\NormalTok{(den\_i))}
\NormalTok{  \}}
  
  \FunctionTok{return}\NormalTok{(}\FunctionTok{list}\NormalTok{(}\AttributeTok{num =}\NormalTok{ num, }\AttributeTok{den =}\NormalTok{ den))}
\NormalTok{\}}
\end{Highlighting}
\end{Shaded}

\subsection{Algorithm: Frequency Response}\label{algorithm-frequency-response}

\begin{Shaded}
\begin{Highlighting}[]
\CommentTok{\# Compute frequency response}
\NormalTok{freq\_response }\OtherTok{\textless{}{-}} \ControlFlowTok{function}\NormalTok{(tf, freq\_range) \{}
\NormalTok{  s }\OtherTok{\textless{}{-}} \DecValTok{1}\DataTypeTok{i} \SpecialCharTok{*}\NormalTok{ freq\_range}
  
\NormalTok{  H }\OtherTok{\textless{}{-}} \FunctionTok{sapply}\NormalTok{(s, }\ControlFlowTok{function}\NormalTok{(si) \{}
\NormalTok{    num\_val }\OtherTok{\textless{}{-}} \FunctionTok{polyval}\NormalTok{(tf}\SpecialCharTok{$}\NormalTok{num, si)}
\NormalTok{    den\_val }\OtherTok{\textless{}{-}} \FunctionTok{polyval}\NormalTok{(tf}\SpecialCharTok{$}\NormalTok{den, si)}
    \FunctionTok{return}\NormalTok{(num\_val }\SpecialCharTok{/}\NormalTok{ den\_val)}
\NormalTok{  \})}
  
\NormalTok{  magnitude\_dB }\OtherTok{\textless{}{-}} \DecValTok{20} \SpecialCharTok{*} \FunctionTok{log10}\NormalTok{(}\FunctionTok{abs}\NormalTok{(H))}
\NormalTok{  phase\_deg }\OtherTok{\textless{}{-}} \FunctionTok{atan2}\NormalTok{(}\FunctionTok{Im}\NormalTok{(H), }\FunctionTok{Re}\NormalTok{(H)) }\SpecialCharTok{*} \DecValTok{180} \SpecialCharTok{/}\NormalTok{ pi}
  
  \FunctionTok{return}\NormalTok{(}\FunctionTok{data.frame}\NormalTok{(}
    \AttributeTok{freq =}\NormalTok{ freq\_range,}
    \AttributeTok{magnitude\_dB =}\NormalTok{ magnitude\_dB,}
    \AttributeTok{phase\_deg =}\NormalTok{ phase\_deg}
\NormalTok{  ))}
\NormalTok{\}}

\CommentTok{\# Helper: polynomial evaluation}
\NormalTok{polyval }\OtherTok{\textless{}{-}} \ControlFlowTok{function}\NormalTok{(p, x) \{}
  \FunctionTok{sum}\NormalTok{(p }\SpecialCharTok{*}\NormalTok{ x}\SpecialCharTok{\^{}}\NormalTok{((}\FunctionTok{length}\NormalTok{(p)}\SpecialCharTok{{-}}\DecValTok{1}\NormalTok{)}\SpecialCharTok{:}\DecValTok{0}\NormalTok{))}
\NormalTok{\}}
\end{Highlighting}
\end{Shaded}

\subsection{Algorithm: Time-Domain Simulation}\label{algorithm-time-domain-simulation}

For network simulation with feedback:

\begin{Shaded}
\begin{Highlighting}[]
\CommentTok{\# Define ODE system}
\NormalTok{network\_ode }\OtherTok{\textless{}{-}} \ControlFlowTok{function}\NormalTok{(t, x, params) \{}
\NormalTok{  N }\OtherTok{\textless{}{-}} \FunctionTok{length}\NormalTok{(x) }\SpecialCharTok{/} \DecValTok{2}  \CommentTok{\# Number of elements}
  
\NormalTok{  dxdt }\OtherTok{\textless{}{-}} \FunctionTok{numeric}\NormalTok{(}\DecValTok{2}\SpecialCharTok{*}\NormalTok{N)}
  
  \ControlFlowTok{for}\NormalTok{ (i }\ControlFlowTok{in} \DecValTok{1}\SpecialCharTok{:}\NormalTok{N) \{}
\NormalTok{    S\_i }\OtherTok{\textless{}{-}}\NormalTok{ x[}\DecValTok{2}\SpecialCharTok{*}\NormalTok{i }\SpecialCharTok{{-}} \DecValTok{1}\NormalTok{]}
\NormalTok{    C\_i }\OtherTok{\textless{}{-}}\NormalTok{ x[}\DecValTok{2}\SpecialCharTok{*}\NormalTok{i]}
    
    \CommentTok{\# Substrate compartment}
\NormalTok{    dxdt[}\DecValTok{2}\SpecialCharTok{*}\NormalTok{i }\SpecialCharTok{{-}} \DecValTok{1}\NormalTok{] }\OtherTok{\textless{}{-}}\NormalTok{ params}\SpecialCharTok{$}\NormalTok{Ra[i] }\SpecialCharTok{{-}}\NormalTok{ params}\SpecialCharTok{$}\NormalTok{k3[i] }\SpecialCharTok{*}\NormalTok{ C\_i}
    
    \CommentTok{\# Controller compartment}
\NormalTok{    dxdt[}\DecValTok{2}\SpecialCharTok{*}\NormalTok{i] }\OtherTok{\textless{}{-}}\NormalTok{ params}\SpecialCharTok{$}\NormalTok{k4[i] }\SpecialCharTok{*}\NormalTok{ S\_i }\SpecialCharTok{{-}}\NormalTok{ params}\SpecialCharTok{$}\NormalTok{k2[i] }\SpecialCharTok{*}\NormalTok{ C\_i}
    
    \CommentTok{\# Add feedback terms from other elements}
    \ControlFlowTok{for}\NormalTok{ (j }\ControlFlowTok{in} \DecValTok{1}\SpecialCharTok{:}\NormalTok{N) \{}
      \ControlFlowTok{if}\NormalTok{ (params}\SpecialCharTok{$}\NormalTok{feedback[i, j] }\SpecialCharTok{!=} \DecValTok{0}\NormalTok{) \{}
\NormalTok{        dxdt[}\DecValTok{2}\SpecialCharTok{*}\NormalTok{i] }\OtherTok{\textless{}{-}}\NormalTok{ dxdt[}\DecValTok{2}\SpecialCharTok{*}\NormalTok{i] }\SpecialCharTok{+}\NormalTok{ params}\SpecialCharTok{$}\NormalTok{feedback[i, j] }\SpecialCharTok{*}\NormalTok{ x[}\DecValTok{2}\SpecialCharTok{*}\NormalTok{j }\SpecialCharTok{{-}} \DecValTok{1}\NormalTok{]}
\NormalTok{      \}}
\NormalTok{    \}}
\NormalTok{  \}}
  
  \FunctionTok{return}\NormalTok{(}\FunctionTok{list}\NormalTok{(dxdt))}
\NormalTok{\}}

\CommentTok{\# Solve}
\FunctionTok{library}\NormalTok{(deSolve)}
\NormalTok{out }\OtherTok{\textless{}{-}} \FunctionTok{ode}\NormalTok{(}\AttributeTok{y =}\NormalTok{ initial\_conditions, }
           \AttributeTok{times =}\NormalTok{ time\_seq, }
           \AttributeTok{func =}\NormalTok{ network\_ode, }
           \AttributeTok{parms =}\NormalTok{ parameters)}
\end{Highlighting}
\end{Shaded}

\chapter*{References}\label{references}
\addcontentsline{toc}{chapter}{References}

\nocite{*}
\bibliography{FRreferences.bib}

\end{document}
